\documentclass{ximera}

\input{../../../preamble.tex}


\definecolor{fvdnavy}{HTML}{071D43}
\definecolor{fvdcyan}{HTML}{20BFC6}
\definecolor{fvdcyanlight}{HTML}{EFFCFB}
\definecolor{fvdmagenta}{HTML}{F10D69}
\definecolor{fvdmagentalight}{HTML}{FFF1F7}
\definecolor{fvdtext}{HTML}{163044}





%=================================================
% Pedagogische omgevingen
% PDF: tcolorbox
% HTML: learning-box
%=================================================

\newenvironment{voorkennis}
{%
    \ifdefined\HCode
        \HCode{%
<section class="learning-box learning-box--prior">
<div class="learning-box__header">
<span class="learning-box__icon" aria-hidden="true"></span>
<span class="learning-box__title">Voorkennis</span>
</div>
<div class="learning-box__content">
        }%
        \begin{itemize}
    \else
       \begin{tcolorbox}[
    enhanced,
    breakable,
    colback=fvdcyanlight,
    colframe=fvdcyan,
    coltext=fvdtext,
    boxrule=0.5pt,
    leftrule=4pt,
    arc=4mm,
    outer arc=4mm,
    left=7mm,
    right=6mm,
    top=5mm,
    bottom=5mm,
    before skip=6mm,
    after skip=6mm,
    title=Voorkennis,
    coltitle=fvdcyan,
    fonttitle=\bfseries\Large,
    attach title to upper={\par\smallskip},
    boxed title style={
        empty,
        left=0mm,
        right=0mm,
        top=0mm,
        bottom=0mm
    },
    overlay={
        \draw[
            fvdcyan,
            line width=0.5pt
        ]
        ([xshift=42mm,yshift=-13mm]frame.north west)
        --
        ([xshift=-7mm,yshift=-13mm]frame.north east);
    }
]
        \begin{itemize}
    \fi
}
{%
    \end{itemize}
    \ifdefined\HCode
        \HCode{%
</div>
</section>
        }%
    \else
        \end{tcolorbox}
    \fi
}


\newenvironment{leerdoelen}
{%
    \ifdefined\HCode
        \HCode{%
<section class="learning-box learning-box--goals">
<div class="learning-box__header">
<span class="learning-box__icon" aria-hidden="true"></span>
<span class="learning-box__title">Leerdoelen</span>
</div>
<div class="learning-box__content">
        }%
        \begin{itemize}
    \else
\begin{tcolorbox}[
    enhanced,
    breakable,
    colback=fvdmagentalight,
    colframe=fvdmagenta,
    coltext=fvdtext,
    boxrule=0.5pt,
    leftrule=4pt,
    arc=4mm,
    outer arc=4mm,
    left=7mm,
    right=6mm,
    top=5mm,
    bottom=5mm,
    before skip=6mm,
    after skip=6mm,
    title=Leerdoelen,
    coltitle=fvdmagenta,
    fonttitle=\bfseries\Large,
    attach title to upper={\par\smallskip},
    boxed title style={
        empty,
        left=0mm,
        right=0mm,
        top=0mm,
        bottom=0mm
    },
    overlay={
        \draw[
            fvdmagenta,
            line width=0.5pt
        ]
        ([xshift=42mm,yshift=-13mm]frame.north west)
        --
        ([xshift=-7mm,yshift=-13mm]frame.north east);
    }
]
        \begin{itemize}
    \fi
}
{%
    \end{itemize}
    \ifdefined\HCode
        \HCode{%
</div>
</section>
        }%
    \else
        \end{tcolorbox}
    \fi
}


\addPrintStyle{../../..}

\title{Oefeningen --- Correlatie}
\author{Ingmar Herreman}



\begin{document}

% FVD_CHAPTER_DATA_START
% Automatisch gegenereerd uit index.tex.
% Niet handmatig aanpassen.
\fvdchapterdata{3}{Verbanden onderzoeken met data}{7}
% FVD_CHAPTER_DATA_END






\maketitle


% ============================================================
\FVDsection{Basis --- r lezen en interpreteren}
% ============================================================

\begin{oefening}[Teken en sterkte]{\oefB \ESS}
Geef voor elke correlatiecoëfficiënt de richting en een globale
beschrijving van de sterkte van het lineaire verband.

\[
\begin{array}{c|ccccc}
& A&B&C&D&E\\ \hline
r&0{,}94&-0{,}88&0{,}12&-0{,}41&1
\end{array}
\]

\begin{oplossing}
\begin{itemize}
  \item $A$: sterk positief lineair;
  \item $B$: sterk negatief lineair;
  \item $C$: zwak positief lineair;
  \item $D$: eerder zwak negatief lineair;
  \item $E$: perfect positief lineair.
\end{itemize}
De woorden ``sterk'' en ``zwak'' zijn globale beschrijvingen; de
context en de puntenwolk blijven belangrijk.
\end{oplossing}
\end{oefening}


\begin{oefening}[Welke is sterker?]{\oefB \ESS}
Vergelijk telkens de sterkte van de lineaire samenhang.

\begin{question}
Welke is sterker?

\[
r_1=-0{,}91,
\qquad
r_2=0{,}84
\]

\begin{multipleChoice}
  \choice[correct]{$r_1$}
  \choice{$r_2$}
  \choice{even sterk}
\end{multipleChoice}
\end{question}

\begin{question}
Waarom?

\begin{oplossing}
Omdat $|-0{,}91|=0{,}91$ groter is dan $|0{,}84|=0{,}84$.
Het minteken geeft alleen de richting aan.
\end{oplossing}
\end{question}

\begin{question}
Welke is sterker?

\[
r_3=-0{,}72,
\qquad
r_4=0{,}72
\]

\begin{multipleChoice}
  \choice{$r_3$}
  \choice{$r_4$}
  \choice[correct]{even sterk}
\end{multipleChoice}
\end{question}
\end{oefening}


\begin{oefening}[Mogelijk of onmogelijk?]{\oefB \ESS}
Duid aan welke waarden een correlatiecoëfficiënt kunnen zijn.

\[
-1{,}4,\quad -1,\quad -0{,}37,\quad 0,\quad 0{,}83,\quad 1,\quad 1{,}02
\]

\begin{oplossing}
Mogelijk zijn

\[
-1,\quad -0{,}37,\quad 0,\quad 0{,}83,\quad 1.
\]

Een correlatiecoëfficiënt ligt steeds tussen $-1$ en $1$.
\end{oplossing}
\end{oefening}


\begin{oefening}[Wat betekent nul?]{\oefB \ESS}
Een dataset heeft

\[
r=0{,}02.
\]

Een leerling besluit:

\begin{quote}
Er bestaat geen verband tussen de twee grootheden.
\end{quote}

Beoordeel deze uitspraak.

\begin{oplossing}
De uitspraak is te sterk. $r=0{,}02$ wijst op vrijwel geen
\emph{lineaire} samenhang. Er kan nog steeds een duidelijk niet-lineair
verband bestaan. Daarom moet ook de puntenwolk bekeken worden.
\end{oplossing}
\end{oefening}


% ============================================================
\FVDsection{Verdieping --- r en puntenwolk samen gebruiken}
% ============================================================

\begin{oefening}[Temperatuur en gasverbruik]{\oefU \ESS}
Voor tien winterdagen werden de gemiddelde buitentemperatuur en het
gasverbruik van een school gemeten.

\[
\begin{array}{c|rrrrrrrrrr}
T\;({}^{\circ}\mathrm C)&-2&0&2&4&5&7&8&10&11&13\\ \hline
G\;(\mathrm{kWh})&820&790&735&690&655&610&590&535&520&470
\end{array}
\]

\begin{question}
Maak met ICT een spreidingsdiagram.
\end{question}

\begin{question}
Bepaal met ICT de correlatiecoëfficiënt $r$.

\begin{oplossing}
Afhankelijk van afronding wordt een waarde zeer dicht bij $-1$
gevonden.
\end{oplossing}
\end{question}

\begin{question}
Leg uit waarom het teken van $r$ past bij de puntenwolk.

\begin{oplossing}
De puntenwolk daalt: hogere buitentemperaturen gaan doorgaans samen
met lager gasverbruik. Daarom is $r$ negatief.
\end{oplossing}
\end{question}

\begin{question}
Formuleer een volledige conclusie.

\begin{oplossing}
De puntenwolk en de correlatiecoëfficiënt wijzen op een zeer sterk
negatief en ongeveer lineair verband tussen buitentemperatuur en
gasverbruik tijdens deze dagen. Hogere temperaturen gaan in deze
dataset doorgaans samen met een lager gasverbruik.
\end{oplossing}
\end{question}
\end{oefening}


\begin{oefening}[Sterk verband, kleine r]{\oefU \ESS}
Een sensor levert de volgende gegevens.

\[
\begin{array}{c|rrrrrrrrr}
x&-4&-3&-2&-1&0&1&2&3&4\\ \hline
y&16&9&4&1&0&1&4&9&16
\end{array}
\]

\begin{question}
Maak een spreidingsdiagram en beschrijf het patroon.

\begin{oplossing}
De punten liggen exact op een U-vormig patroon, overeenkomstig
$y=x^2$. Het verband is zeer sterk, maar niet lineair.
\end{oplossing}
\end{question}

\begin{question}
Bepaal $r$ met ICT.

\begin{oplossing}
Door de symmetrie is $r=0$.
\end{oplossing}
\end{question}

\begin{question}
Waarom is ``$r=0$, dus er is geen verband'' hier fout?

\begin{oplossing}
Omdat $r$ lineaire samenhang meet. De puntenwolk vertoont een perfect
niet-lineair verband.
\end{oplossing}
\end{question}

\begin{question}
Welke algemene regel neem je uit deze oefening mee?

\begin{oplossing}
Bekijk altijd de puntenwolk voordat je een correlatiecoëfficiënt
interpreteert.
\end{oplossing}
\end{question}
\end{oefening}


\begin{oefening}[De invloed van een uitschieter]{\oefU \ESS}
Beschouw eerst de dataset

\[
\begin{array}{c|rrrrrrr}
x&1&2&3&4&5&6&7\\ \hline
y&2&4&6&8&10&12&14
\end{array}
\]

en voeg daarna de waarneming

\[
(8,2)
\]

toe.

\begin{question}
Bepaal $r$ voor de eerste zeven punten.

\begin{oplossing}
$r=1$, want de punten liggen exact op de stijgende rechte $y=2x$.
\end{oplossing}
\end{question}

\begin{question}
Maak daarna de puntenwolk met het extra punt en bepaal opnieuw $r$.
\end{question}

\begin{question}
Beschrijf wat er gebeurt.

\begin{oplossing}
Het extra punt wijkt sterk af van het oorspronkelijke lineaire patroon
en verlaagt de correlatiecoëfficiënt aanzienlijk. Eén uitschieter kan
dus een numerieke samenvatting sterk beïnvloeden.
\end{oplossing}
\end{question}

\begin{question}
Mag je het extra punt daarom verwijderen?

\begin{oplossing}
Niet automatisch. Eerst moet onderzocht worden of het om een fout,
een uitzonderlijke correcte waarneming of een andere situatie gaat.
\end{oplossing}
\end{question}
\end{oefening}


\begin{oefening}[Zelfde richting, andere sterkte]{\oefU}
Twee experimenten leveren beide een positieve correlatie op:

\[
r_A=0{,}96,
\qquad
r_B=0{,}54.
\]

\begin{question}
Wat hebben beide datasets gemeenschappelijk?

\begin{oplossing}
Beide vertonen een positieve lineaire samenhang.
\end{oplossing}
\end{question}

\begin{question}
Wat is verschillend?

\begin{oplossing}
De lineaire samenhang is in dataset $A$ duidelijk sterker dan in
dataset $B$, omdat $|r_A|>|r_B|$.
\end{oplossing}
\end{question}

\begin{question}
Kun je zonder puntenwolken uitsluiten dat één van de datasets een
uitschieter bevat?

\begin{oplossing}
Nee. De correlatiecoëfficiënt alleen toont de structuur van de
puntenwolk niet.
\end{oplossing}
\end{question}
\end{oefening}


% ============================================================
\FVDsection{Geïntegreerd analyseren}
% ============================================================

\begin{oefening}[Absorbantie en concentratie]{\oefU \ESS}
In een chemisch experiment wordt voor verschillende concentraties van
een oplossing de absorbantie gemeten.

\[
\begin{array}{c|rrrrrrrr}
c\;(\mathrm{mmol/L})&0{,}5&1{,}0&1{,}5&2{,}0&2{,}5&3{,}0&3{,}5&4{,}0\\ \hline
A&0{,}12&0{,}21&0{,}34&0{,}43&0{,}55&0{,}64&0{,}77&0{,}85
\end{array}
\]

Voer een volledige correlatieanalyse uit.

\begin{question}
Maak een spreidingsdiagram.
\end{question}

\begin{question}
Beschrijf richting, vorm, sterkte en eventuele bijzonderheden.

\begin{oplossing}
De punten vertonen een sterk positief en ongeveer lineair verband,
zonder opvallende uitschieters.
\end{oplossing}
\end{question}

\begin{question}
Bepaal $r$ met ICT.
\end{question}

\begin{question}
Leg uit hoe de waarde van $r$ aansluit bij de grafische analyse.

\begin{oplossing}
De waarde van $r$ ligt zeer dicht bij $1$, wat past bij de puntenwolk:
de punten liggen dicht rond een stijgende rechte band.
\end{oplossing}
\end{question}

\begin{question}
Formuleer een contextgebonden conclusie.

\begin{oplossing}
Binnen het onderzochte concentratiebereik vertonen concentratie en
absorbantie een zeer sterke positieve en ongeveer lineaire samenhang:
hogere concentraties gaan in deze dataset samen met hogere
absorbanties.
\end{oplossing}
\end{question}
\end{oefening}


\begin{oefening}[Trainingsuren en testresultaat]{\oefU \ESS}
Bij twaalf leerlingen wordt het aantal uren voorbereiding voor een
test vergeleken met het behaalde resultaat. De software geeft

\[
r=0{,}76.
\]

De puntenwolk is stijgend en ongeveer lineair, maar bevat vrij veel
spreiding.

\begin{question}
Interpreteer $r$.

\begin{oplossing}
Er is een vrij sterke positieve lineaire samenhang: leerlingen met
meer voorbereidingsuren behalen in deze dataset doorgaans hogere
resultaten.
\end{oplossing}
\end{question}

\begin{question}
Een leerling schrijft:
``76\% van het testresultaat wordt veroorzaakt door het aantal
trainingsuren.'' Beoordeel die uitspraak.

\begin{oplossing}
Die uitspraak is fout. $r=0{,}76$ is geen percentage verklaarde
variatie en bewijst bovendien geen causaliteit.
\end{oplossing}
\end{question}

\begin{question}
Welke informatie uit de puntenwolk blijft belangrijk naast $r$?

\begin{oplossing}
De ongeveer lineaire vorm, de spreiding van de punten en eventuele
uitschieters of clusters.
\end{oplossing}
\end{question}
\end{oefening}


% ============================================================
\FVDsection{Gevorderd --- wanneer één getal niet genoeg is}
% ============================================================

\begin{oefening}[Maak r misleidend]{\oefG}
Ontwerp met ICT twee datasets:

\begin{enumerate}
  \item een dataset met $r$ dicht bij $0$ maar met een duidelijk
        niet-lineair verband;
  \item een dataset waarin één uitzonderlijk punt de waarde van $r$
        sterk verandert.
\end{enumerate}

Maak voor beide datasets een spreidingsdiagram en schrijf een korte
analyse waarin je uitlegt waarom de puntenwolk noodzakelijk is voor
een correcte interpretatie.

\begin{oplossing}
Er zijn veel correcte oplossingen mogelijk. Bij de eerste dataset
moet een duidelijk krom patroon zichtbaar zijn ondanks een kleine
lineaire correlatie. Bij de tweede moet aantoonbaar zijn dat één
waarneming de correlatiecoëfficiënt sterk beïnvloedt.
\end{oplossing}
\end{oefening}


\begin{oefening}[Kan r alleen beslissen?]{\oefG}
Een onderzoeker heeft drie datasets en wil voor elk ervan later een
lineair model opstellen. Hij sorteert ze uitsluitend volgens $|r|$ en
kiest de dataset met de grootste waarde.

Leg uit waarom deze werkwijze onvoldoende is en stel een betere
beslisprocedure voor.

\begin{oplossing}
Een grote $|r|$ is nuttige informatie over lineaire samenhang, maar
toont niet of er bijvoorbeeld uitschieters, clusters of andere
bijzonderheden aanwezig zijn. Een betere procedure is eerst de
puntenwolk te onderzoeken, de lineariteit en bijzondere punten te
beoordelen, daarna $r$ te interpreteren en pas vervolgens te beslissen
of een lineair model inhoudelijk zinvol is.
\end{oplossing}
\end{oefening}


% ============================================================
\FVDsection{Essentiële oefeningen}
% ============================================================

Voor dit hoofdstuk behoren de volgende oefeningen tot het minimumtraject:

\begin{itemize}
  \item \textbf{Teken en sterkte} --- teken en $|r|$ correct interpreteren;
  \item \textbf{Welke is sterker?} --- sterkte vergelijken via $|r|$;
  \item \textbf{Mogelijk of onmogelijk?} --- het bereik van $r$ beheersen;
  \item \textbf{Wat betekent nul?} --- $r\approx0$ correct interpreteren;
  \item \textbf{Temperatuur en gasverbruik} --- $r$ bepalen en in context analyseren;
  \item \textbf{Sterk verband, kleine r} --- de beperking tot lineaire samenhang begrijpen;
  \item \textbf{De invloed van een uitschieter} --- grafische en numerieke informatie combineren;
  \item \textbf{Absorbantie en concentratie} --- een volledige correlatieanalyse uitvoeren;
  \item \textbf{Trainingsuren en testresultaat} --- correcte en foutieve conclusies onderscheiden.
\end{itemize}

De \oefB-oefeningen leggen de noodzakelijke begrippen en technieken
vast. De geïntegreerde \oefU-oefeningen met \ESS zijn eveneens
essentieel: het leerplandoel vraagt dat leerlingen gegevens
\textbf{analyseren}, niet alleen waarden van $r$ classificeren.

\end{document}
